\textbf{\textit{Soal.}} Diberikan segitiga $ABC$ dengan $\angle CAB = 75\dg$ dan $ABC=15\dg$. Titik $D$ pada segmen $BC$ sehingga $BD=CA$ dan $E$ pada sinar $CA$ sehingga $CE=BC$. Misalkan $F$ titik tengah $AB$ dan $M$ perpotongan $CF$ dengan $BE$. Tunjukkan bahwa $CM$ dan $DM$ menyinggung lingkaran luar segitiga $AEF$.
\begin{proof}[\textbf{Solusi. }] (Sunday, 22 March 2026) 
    \begin{center}
\definecolor{wewdxt}{rgb}{0.43137254901960786,0.42745098039215684,0.45098039215686275}
\definecolor{rvwvcq}{rgb}{0.08235294117647059,0.396078431372549,0.7529411764705882}
\begin{tikzpicture}[line cap=round,line join=round,>=triangle 45,scale=3]
\clip(-0.2449028041008735,-0.6003285076290579) rectangle (3.6,3.2);
\draw [shift={(0.,3.)},line width=1pt,fill=black,fill opacity=0.0100000001490112] (0,0) -- (-90.:0.3966815960494947) arc (-90.:-75.:0.3966815960494947) -- cycle;
\draw [line width=0.7pt,color=wewdxt] (0.,0.)-- (0.,3.);
\draw [line width=0.7pt,color=wewdxt] (0.,3.)-- (0.803847577293368,0.);
\draw [line width=0.7pt,color=wewdxt] (0.803847577293368,0.)-- (0.,0.);
\draw [line width=0.7pt,color=wewdxt] (0.,0.) circle (3.cm);
\draw [line width=0.7pt,color=wewdxt] (0.803847577293368,0.)-- (3.,0.);
\draw [line width=0.7pt,color=wewdxt] (0.6339745962155613,2.366025403784439)-- (0.,0.);
\draw [line width=0.7pt,color=wewdxt] (1.9019237886466838,1.098076211353316) circle (1.5529142706151244cm);
\draw [line width=0.7pt,color=wewdxt] (0.,2.196152422706632)-- (1.5,2.5980762113533156);
\draw [line width=0.7pt,color=wewdxt] (0.,3.)-- (3.,0.);
\draw [line width=0.7pt,color=wewdxt] (0.,3.)-- (1.5,2.5980762113533156);
\draw [line width=0.7pt,color=wewdxt] (0.,0.)-- (1.5,2.5980762113533156);
\begin{scriptsize}
\draw [fill=rvwvcq] (0.,0.) circle (1pt);
\draw[color=rvwvcq] (0.0,-0.1) node {$C$};
\draw [fill=rvwvcq] (0.,3.) circle (1pt);
\draw[color=rvwvcq] (0.09888791247535525,3.148312575038661) node {$B$};
\draw[color=black] (0.061,2.5) node {$15\dg$};
\draw [fill=wewdxt] (0.803847577293368,0.) circle (1pt);
\draw[color=wewdxt] (0.8922511045743448,0.21286876427240492) node {$A$};
\draw [fill=wewdxt] (0.,2.196152422706632) circle (1pt);
\draw[color=wewdxt] (0.09888791247535525,2.3) node {$D$};
\draw [fill=wewdxt] (3.,0.) circle (1pt);
\draw[color=wewdxt] (3.087222602714882,0.21286876427240492) node {$E$};
\draw [fill=wewdxt] (0.401923788646684,1.5) circle (1pt);
\draw[color=wewdxt] (0.49556950852484993,1.5) node {$F$};
\draw [fill=wewdxt] (0.6339745962155613,2.366025403784439) circle (1pt);
\draw[color=wewdxt] (0.7203557462862303,2.5797356207010527) node {$M$};
\draw [fill=wewdxt] (1.5,2.5980762113533156) circle (1pt);
\draw[color=wewdxt] (1.5930552575951187,2.8177445783307493) node {$G$};
\end{scriptsize}
\end{tikzpicture}
\end{center}

    Misalkan $G$ titik pada sinar $DM$ sedemikian sehingga $MG = MF$ dan $M$ berada di antara $D$ dan $G$.

\textbf{Bukti $CF$ menyinggung $(AEF)$ di $F$}\\
Perhatikan bahwa $\triangle ABC$ siku-siku di $C$. Karena $F$ titik tengah $AB$ maka jelas bahwa $BF=AF=CF$. Selain itu, jelas bahwa $\angle FAC = \angle ACF = 75\dg$ sehingga $\angle CFA = 30\dg$ dan $\angle BFC = 150\dg$ dengan $\sin CFA = \sin BFC = \frac{1}{2}$. Oleh karena itu, dengan ditambah fakta $CB = CE$ didapat
\begin{align*}
    [ABC] &= [CFA] + [CFB]\\
    \frac{1}{2} \times CA \times CB &= \frac{1}{2} \times AF \times FC \times \sin \angle CFA + \frac{1}{2} \times BF \times FC \times \sin \angle BFC\\
     CA \times CB &= AF \times FC \times \frac{1}{2} + BF \times FC \times \frac{1}{2} \\
     CA \times CE &= CF^2
\end{align*}
yang menunjukkan bahwa $CF$ menyinggung $(AEF)$ di $F$ dari \textit{Power of a Point}.\\
\newline
\textbf{Bukti $MG$ menyinggung $(AEF)$ di $G$}\\
Dikarenakan $CB = CE \implies \angle DBM = \angle MEC = 45\dg$ dan $\angle MCE = \angle FCA = 75\dg$ maka $\angle CME = 60\dg$.
Lalu, karena $BD = CA$, $BC = CE$, dan $AF=FB$, dengan Teorema Menelaus pada segitga $EMC$ dan garis transversal $AB$ diperoleh
\begin{align*}
    \dfrac{AF}{FB} \times \dfrac{BM}{ME} \times \dfrac{EC}{CA} &= 1\\
    1 \times \dfrac{BM}{ME} \times \dfrac{EC}{BD} &= 1\\
    \dfrac{BM}{ME} &= \dfrac{BD}{EC}.
\end{align*}
dan $\angle DBM = \angle MEC$ maka $\triangle BDM \sim \triangle ECM$. Sehingga didapat $\angle EMG = \angle BMD = \angle CME = 60\dg$. 

Oleh karena itu $MF = MG$, $\angle EMG = \angle CME$, dan fakta bahwa $MF$ menyinggung $(AEF)$ di $F$, akan mengakibatkan $MG$ menyinggung $(AEF)$ di $G$.

Dari kedua kasus tersebut, terbukti bahwa $CM$ dan $DM$ menyinggung lingkaran luar segitiga $AEF$.
\end{proof}